\documentclass[11pt]{article}
\usepackage[T1]{fontenc}
\usepackage[utf8]{inputenc}
\usepackage{mathptmx,courier}
\usepackage[margin=2.7cm]{geometry}
\usepackage{amsmath,amssymb,booktabs,array,graphicx,float}
\usepackage{setspace,microtype}
\usepackage[numbers,sort&compress]{natbib}
\usepackage[colorlinks=true,linkcolor=blue,citecolor=blue,urlcolor=blue]{hyperref}
\setstretch{1.18}
\raggedbottom
\AtBeginDocument{\lefthyphenmin=2\righthyphenmin=3}
\renewcommand{\topfraction}{0.9}
\renewcommand{\textfraction}{0.05}
\renewcommand{\floatpagefraction}{0.85}
\setlength{\emergencystretch}{2em}
\newcommand{\ind}{\mathbf{1}}
\newcommand{\equalcontribution}{Siyuan Tong and Yuechen Wang contributed equally to this work and share first authorship.}
\newcommand{\jointauthors}{\begin{minipage}{0.96\textwidth}\centering Siyuan Tong$^\dagger$ and Yuechen Wang$^\dagger$\\[0.3em]\small $^\dagger$\equalcontribution\end{minipage}}
\newcommand{\papertitle}{Separating coverage gains from error targeting in molecular prediction sets}
\hypersetup{pdftitle={Separating coverage gains from error targeting in molecular prediction sets},pdfauthor={Siyuan Tong and Yuechen Wang}}

\begin{document}
\noindent 12 November 2026

\noindent Dear Dr. Park,

Following your decision on \textbf{CHEMOLAB-D-26-01118}, we submit a substantially revised version as a \textbf{new original research article} in \emph{Chemometrics and Intelligent Laboratory Systems}. The previous title was \emph{Beyond Aggregate Reliability in ADMET Prediction: Protecting Weakest-Class Coverage under Molecular Distribution Shift with a Transport-Audited Multi-View Envelope}. The revised title is \emph{\papertitle}.

The revised study introduces class-budget matching to isolate the error-targeting value of molecular prediction-set expansion and an exact weighted-rank condition connecting transport weights to label addition. In response to the reviewer's comments, we have made the following substantive changes.

\begin{enumerate}
\item \textbf{Contribution and the earlier Stage I examples.} The distinction between marginal and class-conditional coverage now serves as established background. The contribution centers on matching expansion counts within each original singleton class, decomposing recovered errors and withheld correct decisions, and characterizing activation. The new experiments use the corrected finite-sample quantile rule, with an infinite threshold when the required rank exceeds the calibration sample size. The earlier Stage I numerical examples and figures have been replaced by the newly executed evidence.

\item \textbf{Set expansion and transport gates.} We distinguish view eligibility, actual label addition, and error recovery. The exact activation inequality and constant-query-weight ranking identity explain how the weighted rule operates. In the primary active-label comparison, both views pass in all 30 assay--seed cells but produce no additional deferrals; these results are now explained through the activation mechanism.

\item \textbf{Chemical decision objective.} Historical label-0 protection and the prespecified label-1 objective are reported separately. The revised text defines the measured assay outcomes, false-inactive and false-active exposure, correct-singleton yield, and operational deferral. It also accounts for missing assay outcomes through their corresponding labeled denominators.

\item \textbf{Comparator strength and robustness.} The completed evaluation contains 960 predictor units across twelve Tox21 assays and fixed random/scaffold holdouts from four additional assay families. It adds shared-predictor, class-budget comparisons; complete CoDrug and KMM-CP reference systems; predictor and score sensitivities; component and gate ablations; and paired molecular bootstrap intervals. Omitted-label probability ranking reproduces the historical 1.3649-point coverage gain at the same class-specific workload and achieves higher observed coverage and yield than generic uncertainty. The revised paper reports the gain together with its deferral cost and the corresponding paired contrast.

\newpage
\item \textbf{Algorithm presentation and reproducibility.} The methods now proceed from cohort roles and set construction to activation, matched budgets, decision metrics, and inference. An activation table replaces repeated zero-difference graphics, and the supporting information specifies gate settings and reference configurations. The submission includes the full numerical supplement, frozen-input rebuilding scripts, and the verified executed CoDrug and KMM-CP sources. The accompanying response maps each review topic to the revised sections, tables, and figures.
\end{enumerate}

These changes provide a practical framework for assessing the information gained from molecular uncertainty procedures at a specified decision workload. We believe this contribution is relevant to the journal's readership in chemometric prediction and uncertainty assessment.

\equalcontribution\ Siyuan Tong remains the corresponding author.

Thank you for considering this substantially revised new submission.

\noindent Sincerely,\\
Siyuan Tong\\
On behalf of the authors\\
\texttt{25064241@siswa.um.edu.my}
\end{document}
